\begin{textsample}{POS Dim 4 – 2021\_c4 – Score 2.00 – 2021\_c4/enade\_2021\_ci\_comp\_9c8a31e088.txt}  \label{ex:f4_pos_136}
CAPÍTULO 4 PERCEPÇÃO DA PROVA Nas análises feitas neste capítulo, constam as percepções dos concluintes da Área de Ciência da Computação ( Licenciatura e Bacharelado ) sobre a prova aplicada no Enade/2021.
Essas percepções foram mensuradas por meio de nove questões que avaliaram desde o grau de dificuldade da prova até o tempo gasto para concluí -la.
As percepções sobre a prova foram relacionadas ao desempenho dos estudantes, à Grande Região de funcionamento do curso, à Categoria Administrativa e ao tipo de Organização Acadêmica da IES.
O questionário de percepção da prova encontra -se ao final dos Anexos VIIIA ( Licenciatura ) e VIIIB ( Bacharelado ), onde está a reprodução do exame.
O desempenho dos estudantes foi classificado em quatro quartos.
Para tanto, esse desempenho foi ordenado de forma ascendente.
O percentil 25, P25, também conhecido como primeiro quartil, é a nota de desempenho que deixa um quarto ( 25 \% ) dos valores observados abaixo e três quartos acima.
Na Figura 1, apresenta -se uma ilustração desse conceito.
O quarto inferior de desempenho é composto pelas notas abaixo do primeiro quartil.
Já o percentil 75, P75, também conhecido como terceiro quartil, é o valor para o qual há três quartos ( 75 \% ) dos dados abaixo e um quarto acima dele.
O quarto superior de desempenho é composto pelas notas iguais ou acima do terceiro quartil.
O percentil 50, P50, também conhecido como mediana, é o valor que divide as notas em dois conjuntos de igual tamanho.
O segundo quarto inclui valores entre o primeiro quartil ( P25 ) e a mediana.
O terceiro quarto contém os valores entre a mediana ( P50 ) e o terceiro quartil ( P75 ).
Vale ressaltar que percentis, quartis e medianas são pontos que, não obrigatoriamente, pertencem ao conjunto original de dados, ao passo que os quartos são subconjuntos dos dados originais.
A seguir, serão apresentados gráficos com resultados selecionados, relativos às nove questões avaliadas por grupos de estudantes.
Nas barras dos gráficos, apresenta -se o percentual de alunos que assinalaram uma das opções ou a soma das porcentagens daqueles que assinalaram duas ( ou três ) delas.
Por exemplo, para as questões 1 e 2, nos gráficos, é apresentada a porcentagem total de participantes que assinalaram as opções ( D ) Difícil e ( E ) Muito difícil.
Em cada barra, foram assinalados, também, os extremos do intervalo de confiança de 95 \% como linhas verticais unidas por uma linha horizontal na forma da letra H maiúscula.
O estimador de um parâmetro com um certo nível de confiança ( por exemplo, 95,0 \% ) deve conter o parâmetro no intervalo de confiança em 95 \% das vezes.
Na comparação entre os estimadores dos parâmetros de duas classes de uma dada categoria ( por exemplo, Norte e Nordeste nas Grandes Regiões, de primeiro e último quarto dentro de desempenho, de IES Públicas e Privadas, ou de Universidades e Faculdades ), associados aos seus respectivos intervalos de confiança, diz -se que não há diferença estatisticamente significativa entre os parâmetros das duas categorias, quando há interseção entre os intervalos de confiança ; e que há diferença, se os intervalos de confiança são disjuntos ( para maiores informações vide Glossário ).
Nas Tabelas do Anexo IIA ( Licenciatura ) e do Anexo IIB ( Bacharelado ) são apresentados os valores absolutos e a distribuição percentual¹⁸ das alternativas válidas das nove questões, segundo o mesmo recorte de desempenho dos alunos e da Grande Região de funcionamento do curso, de Categoria Administrativa e do tipo de Organização Acadêmica da IES. 4.1 GRAU DE DIFICULDADE DA PROVA 4.1.1 COMPONENTE DE FORMAÇÃO GERAL Ao avaliarem “ Qual o grau de dificuldade desta prova na parte de Formação Geral? ” ( Questão 1 ), 43,7 \% do grupo de inscritos e presentes de Licenciatura optaram pelas alternativas Difícil ou Muito difícil, entretanto, para cerca de metade dos estudantes ( 49,7 \% ), o Componente de Formação Geral da prova foi avaliado, quanto ao grau de dificuldade, como Médio ( Gráfico 4.1a, Gráfico 4.2a e, no Anexo IIA, a Tabela II.1a ).
O percentual de estudantes de Licenciatura que consideraram a prova Difícil ou Muito difícil foi maior na região Norte, onde a proporção foi de 47,8 \%, enquanto a de menor proporção foi na Centro-Oeste, com 29,8 \%.
No Gráfico 4.1a, é possível observar que as diferenças entre a região Centro-Oeste e as regiões Norte e Nordeste são estatisticamente significativas.
Nas Grandes Regiões, a proporção de presentes à prova que avaliaram o Componente de Formação Geral, quanto ao grau de dificuldade, como Médio esteve entre 45,6 \%, na região Norte, e 56,0 \%, na região Centro-Oeste.
Ao avaliarem “ Qual o grau de dificuldade desta prova na parte de Formação Geral? ” ( Questão 1 ), 22,5 \% do grupo de inscritos e presentes de Bacharelado optaram pelas alternativas Difícil ou Muito difícil, entretanto, para mais da metade dos estudantes de Bacharelado ( 61,3 \% ), o Componente de Formação Geral da prova, quanto ao grau de dificuldade, foi considerado como Médio ( Gráfico 4.1b, Gráfico 4.2b e, no Anexo IIB, a Tabela II.1b ).
O percentual de estudantes de Bacharelado que consideraram a prova como Difícil ou Muito difícil foi maior na região Norte, onde a proporção foi de 26,3 \%, enquanto a de menor incidência foi a Nordeste, com 18,5 \%.
No Gráfico 4.1b, é possível observar que apenas as diferenças entre a região Nordeste e as regiões Norte, Sudeste e Centro-Oeste são estatisticamente significativas.
Nas Grandes Regiões, a proporção de presentes à prova que consideraram o Componente de Formação Geral como sendo de grau de dificuldade Médio esteve entre 59,6 \%, na região Sudeste, e 63,8 \%, na região Nordeste.
O percentual de alunos de Licenciatura que consideraram o componente de Formação Geral da prova Difícil ou Muito difícil foi decrescente em função dos quartos de desempenho : 54,2 \%, no primeiro quarto, e 30,5 \%, no quarto, o grupo de melhor desempenho na prova.
Nos quartos de desempenho intermediários, a proporção de alunos que consideraram a prova Difícil ou Muito difícil foi de 48,1 \%, no segundo quarto, e de 41,8 \%, no terceiro quarto.
As diferenças entre o quarto superior de desempenho e os demais é estatisticamente significativa, assim como entre o primeiro e o terceiro quartos de desempenho.
Para quase todos os quartos de desempenho, a alternativa modal para essa pergunta foi a Médio, com 47,0 \%, 54,0 \% e 57,0 \% dos respondentes no segundo, no terceiro e no último quarto de desempenho, respectivamente.
No quarto de desempenho inferior, a alternativa modal para essa pergunta foi a Difícil, com 42,1 \% das respostas, enquanto a alternativa Médio alcançou 40,7 \% das respostas.
O percentual de alunos de Bacharelado que consideraram o componente de Formação Geral da prova Difícil ou Muito difícil foi decrescente em função dos quartos de desempenho : 33,0 \%, no primeiro quarto, e 16,1 \%, no quarto, o grupo de melhor desempenho na prova.
Nos quartos de desempenho intermediários, a proporção de alunos que consideraram a prova Difícil ou Muito difícil foi de 23,0 \%, no segundo quarto, e de 18,0 \%, no terceiro quarto.
Apenas a diferença entre os dois quartos superiores de desempenho não é estatisticamente significativa.
Para todos os quartos de desempenho, a alternativa modal para essa pergunta foi a Médio, com 55,2 \% e 63,9 \% dos respondentes nos quartos extremos, primeiro e quarto, respectivamente.
O percentual de estudantes de Licenciatura que consideraram o componente de Formação Geral da prova Difícil ou Muito difícil foi maior entre os estudantes de Faculdades, alcançando 50,0 \%, enquanto entre os estudantes dos Centros Universitários, a Organização Acadêmica com a menor incidência, alcançou 40,8 \%.
No Gráfico 4.3a, é possível observar que não há diferenças estatisticamente significativas entre as Organizações Acadêmicas.
Nas Organizações Acadêmicas, a proporção de presentes de Licenciatura à prova que avaliaram o Componente de Formação Geral como sendo de grau de dificuldade Médio esteve entre 48,5 \%, nos CEFET/IF, e 51,2 \%, nas Universidades ( ver também a Tabela II.2a, no Anexo IIA ).
O percentual de estudantes de Bacharelado que consideraram o componente de Formação Geral da prova como Difícil ou Muito difícil foi maior entre os estudantes de Faculdades, para os quais a proporção foi de 30,7 \%, enquanto entre os de Universidades, a Organização Acadêmica em que foi registrada a menor incidência, foi de 21,3 \%.
No Gráfico 4.3b, é possível observar que as diferenças entre Universidades e Faculdades são estatisticamente significativas, assim como entre Centros Universitários e Faculdades.
Nas Organizações Acadêmicas, a proporção de presentes de Bacharelado à prova que consideraram o Componente de Formação Geral como sendo de grau de dificuldade Médio esteve entre 55,1 \%, nas Faculdades, e 62,3 \%, nos Centros Universitários ( Ver também a Tabela II.2b, no Anexo IIB ).
O percentual de estudantes de Licenciatura que consideraram o componente de Formação Geral da prova Difícil ou Muito difícil foi maior entre os estudantes de cursos de IES Públicas ( 44,0 \% ), sem diferença estatisticamente significativa da proporção registrada nos de IES Privadas ( 40,8 \% ).
Já a proporção de presentes de Licenciatura à prova que avaliaram o Componente de Formação Geral, quanto ao de grau de dificuldade, como Médio foi de 49,8 \%, entre os de IES Públicas, e de 48,3 \%, entre os de IES Privadas ( ver Gráfico 4.4a ).
O percentual de estudantes de Bacharelado que consideraram o componente de Formação Geral da prova como Difícil ou Muito difícil foi maior entre os estudantes de cursos de IES Privadas ( 25,5 \% ), uma diferença estatisticamente significativa da proporção registrada em IES Públicas ( 19,0 \% ).
Já a proporção de presentes de Bacharelado à prova que consideraram o Componente de Formação Geral, quanto ao grau de dificuldade, como Médio foi de 63,4 \%, entre os de IES Públicas, e de 59,6 \%, entre os de IES Privadas ( ver Gráfico 4.4b ). 4.1.2 COMPONENTE DE CONHECIMENTO ESPECÍFICO Ao responderem à Questão 2 – “ Qual o grau de dificuldade desta prova na parte de Componente Específico? ” – 64,8 \% do grupo de estudantes de Licenciatura classificaram -na como Difícil ou Muito difícil.
Além disso, o Componente modal de Conhecimento Específico da prova foi considerado como de grau de dificuldade Médio por 33,0 \% dos alunos ( Gráfico 4.5a, Gráfico 4.6a, e, no Anexo IIA, a Tabela II.3a ).
A análise das respostas dos estudantes de Licenciatura, quanto ao grau de dificuldade do Componente de Conhecimento Específico da prova, agregado por Grande Região, mostra diferenças estatisticamente significativas entre proporções de alunos que a avaliaram como Difícil ou Muito difícil.
Tais diferenças são observadas entre as regiões Norte ( 56,9 \% ) e Sul ( 70,5 \% ), e entre a região Centro-Oeste ( 44,0 \% ) e as regiões Nordeste ( 68,0 \% ), Sudeste ( 66,1 \% ) e Sul ( 70,5 \% ).
O percentual de alunos que classificaram o grau de dificuldade como Médio, no Componente de Conhecimento Específico, variou de 26,8 \% a 52,4 \%, nas regiões Sul e Centro-Oeste, respectivamente.
Ao responderem à Questão 2 – “ Qual o grau de dificuldade desta prova na parte de Componente Específico? ” – 49,9 \% do grupo de estudantes de Bacharelado classificaram -na como Difícil ou Muito difícil.
Além disso, o Componente modal de Conhecimento Específico da prova foi considerado, quanto ao grau de dificuldade, como Médio por 46,4 \% dos alunos ( Gráfico 4.5b, Gráfico 4.6b, e, no Anexo IIB, a Tabela II.3b ).
A análise das respostas dos estudantes de Bacharelado quanto ao grau de dificuldade do Componente de Conhecimento Específico da prova, agregado por Grande Região, mostra que há diferenças estatisticamente significativas entre as proporções de alunos que a avaliaram como Difícil ou Muito difícil.
Tais diferenças são observadas entre as regiões Norte ( 56,0 \% ) e Nordeste ( 44,7 \% ), e entre as regiões Nordeste ( 44,7 \% ) e Sudeste ( 52,0 \% ).
O percentual de alunos que classificaram o grau de dificuldade como Médio, no Componente de Conhecimento Específico, variou de 41,0 \% a 50,6 \%, nas regiões Norte e Nordeste, respectivamente.
Considerando -se a avaliação da dificuldade das questões do Componente de Conhecimento Específico da prova, de acordo com o desempenho dos estudantes de Licenciatura, observa -se uma tendência crescente dos que consideraram a parte específica Difícil ou Muito difícil, variando de 63,8 \% ( primeiro quarto ) a 65,8 \% ( último quarto ).
As demais proporções foram de 64,8 \% e 64,9 \% no segundo e terceiro quartos, respectivamente.
Observa -se que não há diferença estatisticamente significativa dos resultados entre os quartos de desempenho.
Já os que responderam que o grau de dificuldade das questões do Componente de Conhecimento Específico da prova foi Médio, atingiram percentagens que variaram de 34,5 \%, no primeiro quarto, a 31,1 \%, no último.
Considerando -se a avaliação da dificuldade das questões do Componente de Conhecimento Específico da prova, de acordo com o desempenho dos estudantes de Bacharelado, observa -se que só há diferença estatisticamente significativa dos resultados entre o quarto superior de desempenho e os demais quartos.
A proporção dos que classificaram a parte específica como Difícil ou Muito difícil variou de 40,3 \% ( último quarto ) a 54,0 \% ( segundo quarto ).
As demais proporções dessa \textbf{reposta} foram de 53,5 \% e 51,7 \%, nos quartos inferior e terceiro, respectivamente.
Já a proporção dos que responderam que o grau de dificuldade das questões do Componente de Conhecimento Específico da prova foi Médio variou de 41,8 \%, no primeiro quarto, a 55,1 \% no último.
Considerando -se a avaliação da dificuldade das questões do Componente de Conhecimento Específico da prova, o percentual de estudantes de Licenciatura que as consideraram Difícil ou Muito difícil foi maior entre os estudantes de CEFET/IF, atingindo 65,9 \% ; já entre os de Faculdades, a Organização Acadêmica com a menor incidência, essa proporção foi de 50,0 \%.
No Gráfico 4.7a, é possível observar que não há diferença estatisticamente significativa entre as Organizações Acadêmicas.
Nas Organizações Acadêmicas, a proporção de presentes de Licenciatura à prova que consideraram o Componente de Conhecimento Específico de grau de dificuldade Médio esteve entre 31,7 \%, nos Centros Universitários, e 50,0 \%, nas Faculdades ( ver também Tabela II.4a, no Anexo IIA ).
Considerando -se a avaliação da dificuldade das questões do Componente de Conhecimento Específico da prova, o percentual de estudantes de Bacharelado que as consideraram como Difícil ou Muito difícil foi maior entre os de CEFET/IF, tendo alcançado 60,4 \%, enquanto entre os de Universidades, a Organização Acadêmica com a menor incidência, alcançou 47,8 \%.
No Gráfico 4.7b, é possível observar que as diferenças entre Universidades e as demais Organizações Acadêmicas são estatisticamente significativas.
Nas Organizações Acadêmicas, a proporção de presentes de Bacharelado à prova que consideraram o Componente de Formação Geral como de grau de dificuldade Médio esteve entre 36,1 \%, nos CEFET/IF e 48,3 \% nas Universidades ( ver também Tabela II.4b, no Anexo IIB ).
O percentual de estudantes de Licenciatura que consideraram o Componente Específico da prova Difícil ou Muito difícil foi maior entre os estudantes de cursos de IES Públicas ( 65,0 \% ), sem diferença estatisticamente significativa da proporção registrada em IES Privadas ( 63,9 \% ).
Já a proporção de presentes de Licenciatura à prova que avaliaram esse componente como sendo de grau de dificuldade Médio foi de 33,1 \% entre os de IES Públicas e de 32,0 \% entre os de IES Privadas ( ver Gráfico 4.8a e Tabela II.4a, no Anexo IIA ).
O percentual de estudantes de Bacharelado que consideraram o Componente Específico da prova como Difícil ou Muito difícil foi maior entre os estudantes de cursos de IES Privadas ( 54,5 \% ), uma diferença estatisticamente significativa da proporção alcançada em IES Públicas ( 44,5 \% ).
Já a proporção de presentes de Bacharelado à prova que consideraram esse componente como sendo de grau de dificuldade Médio foi de 50,9 \%, entre os de IES Públicas, e de 42,5 \%, entre os de IES Privadas ( ver Gráfico 4.8b e Tabela II.4b, no Anexo IIB ). 4.2 EXTENSÃO DA PROVA EM RELAÇÃO AO TEMPO TOTAL Indagados quanto à extensão da prova, em relação ao tempo total oferecido para a sua resolução ( Questão 3 ), os estudantes de Licenciatura apontaram, com maior incidência, a alternativa que considerava a extensão Adequada, para todas as agregações consideradas ( Gráfico 4.9a, Gráfico 4.10a, e, no Anexo IIA, a Tabela II.5a ).
O percentual de alunos de Licenciatura que responderam ser a extensão da prova Adequada foi de 56,2 \%.
Já 37,2 \% dos inscritos presentes consideraram que a prova foi Longa ou Muito longa, e 6,5 \% a avaliaram como Curta ou Muito curta.
Entre as Grandes Regiões, a avaliação de que a prova foi Longa ou Muito longa, em relação ao tempo total destinado à sua resolução, atingiu porcentagens que variaram de 31,1 \%, na região Norte, até 40,9 \%, na região Sudeste.
Não há diferença estatisticamente significativa entre as Grandes Regiões.
Indagados quanto à “ extensão da prova, em relação ao tempo total ” oferecido para a sua resolução ( Questão 3 ), os estudantes de Bacharelado apontaram, com maior incidência, a alternativa em que se considerava a extensão Adequada, em todas as agregações consideradas ( Gráfico 4.9b, Gráfico 4.10b, e, no Anexo IIB, a Tabela II.5b ).
O percentual de alunos de Bacharelado que responderam ser a extensão da prova Adequada foi de 59,4 \%.
Já 33,0 \% dos inscritos presentes consideraram que a prova foi Longa ou Muito longa, e 7,6 \% a avaliaram como Curta ou Muito curta.
Em relação às Grandes Regiões, a proporção daqueles que avaliaram a prova como Longa ou Muito longa, em relação ao tempo total destinado à sua resolução, variou de 28,4 \%, na região Centro-Oeste, até 35,3 \%, na região Sul.
Há diferença estatisticamente significativa entre a região Centro-Oeste e as regiões Nordeste e Sul.
Considerando -se o desempenho dos alunos de Licenciatura, nota -se ainda que 55,8 \% consideraram a extensão da prova Adequada, no quarto de desempenho inferior, e 56,9 \%, no de melhor desempenho ( quarto superior ).
Nos quartos intermediários, essa proporção foi de 61,0 \%, no segundo quarto, e de 51,4 \%, no terceiro.
No Gráfico 4.10a, observa -se que as proporções de estudantes de Licenciatura que consideraram a prova Longa ou Muito longa não apresentam tendência específica em função dos quartos de desempenho, sendo a maior proporção dessa resposta no terceiro quarto de desempenho ( 40,2 \% ), e a menor, no segundo quarto ( 33,5 \% ).
Pode -se constatar que não há diferença estatisticamente significativa entre as proporções dos quartos de desempenho.
Considerando -se o desempenho dos alunos de Bacharelado, nota -se ainda que 60,6 \% consideraram a extensão da prova Adequada, no quarto de desempenho inferior, e 57,1 \%, no de melhor desempenho ( quarto superior ).
Nos quartos intermediários, essa proporção foi de 62,2 \%, no segundo, e de 57,4 \%, no terceiro quarto.
No Gráfico 4.10b, pode -se constatar que há diferença estatisticamente significativa entre as proporções de estudantes de Bacharelado que consideraram a prova Longa ou Muito longa em função dos quartos de desempenho : entre o primeiro quarto ( 31,2 \% ) e o quarto superior ( 36,0 \% ).
Considerando -se a extensão da prova, o percentual de estudantes de Licenciatura que avaliaram a prova como Longa ou Muito longa, em relação ao tempo total destinado à sua resolução, foi maior entre os estudantes de Faculdades, alcançando 100,0 \%, enquanto entre os estudantes de Universidades, a Organização Acadêmica com a menor incidência, esse índice atingiu 36,2 \%.
No Gráfico 4.11a, é possível observar que há diferença estatisticamente significativa apenas entre as Faculdades e as demais Organizações Acadêmicas.
A proporção de presentes de Licenciatura à prova que consideraram a sua extensão Adequada, em relação ao tempo total destinado à sua resolução, esteve entre 0,0 \%, nas Faculdades, e 57,3 \%, nas Universidades.
Ver também Tabela II.6a, no Anexo IIA.
Considerando -se a extensão da prova, o percentual de estudantes de Bacharelado que avaliaram a prova como Longa ou Muito longa, em relação ao tempo total destinado à sua resolução, foi maior entre os estudantes de Universidades, para os quais a proporção foi de 33,5 \%, enquanto os de Faculdades, a Organização Acadêmica com a menor incidência, foi de 28,0 \%.
No Gráfico 4.11b, é possível observar que somente a diferença entre Universidades e Faculdades é estatisticamente significativa.
Nas Organizações Acadêmicas, a proporção de presentes de Bacharelado à prova que consideraram a sua extensão como Adequada, em relação ao tempo total destinado à sua resolução, esteve entre 57,4 \%, nos CEFET/IF, e 64,1 \%, nas Faculdades ( Ver também Tabela II.6b, no Anexo IIB ).
O percentual de estudantes de Licenciatura que consideraram a prova Longa ou Muito longa, em relação ao tempo total destinado à sua resolução, foi maior entre os estudantes de cursos de IES Privadas ( 42,2 \% ), não havendo diferença estatisticamente significativa da proporção em relação aos de IES Públicas ( 36,7 \% ).
Já a proporção de presentes de Licenciatura à prova que consideraram a extensão Adequada foi de 57,1 \%, entre os de IES Públicas, e de 49,0 \%, entre os de IES Privadas ( ver Gráfico 4.12a e Tabela II.6a, no Anexo IIA ).
O percentual de estudantes de Bacharelado que consideraram a prova como Longa ou Muito longa, em relação ao tempo total destinado à sua resolução, foi maior entre os estudantes de cursos de IES Públicas ( 34,9 \% ), com diferença estatisticamente significativa da proporção registrada em IES Privadas ( 31,5 \% ).
Já a proporção de presentes de Bacharelado à prova que consideraram a extensão como Adequada foi de 57,2 \%, entre os de IES Públicas, e de 61,2 \%, entre os de IES Privadas ( ver Gráfico 4.12b e Tabela II.6b, no Anexo IIB ). 4.3 COMPREENSÃO DOS ENUNCIADOS DAS QUESTÕES 4.3.1 COMPONENTE DE FORMAÇÃO GERAL Relativamente aos enunciados das questões do Componente de Formação Geral ( Questão 4 ), as opiniões foram positivas, já que 71,2 \% dos alunos de Licenciatura avaliados consideraram todos ou a maioria dos enunciados das questões “ claros e objetivos ” ( Gráfico 4.13a, Gráfico 4.14a e, no Anexo IIA, a Tabela II.7a ).
Na análise regional, a percentagem de estudantes de Licenciatura que avaliaram que todos ou a maioria dos enunciados das questões do Componente de Formação Geral estavam “ claros e objetivos ” variou de 67,4 \%, na região Norte, a 85,7 \%, na região Centro-Oeste.
A análise das percepções dos estudantes de Licenciatura sobre a clareza e a objetividade dos enunciados permite afirmar que todos ou a maioria dos enunciados de questões relativas ao Componente de Formação Geral, foram considerados “ claros e objetivos ” para a maior parte dos respondentes ( maior ou igual a 67,4 \%, em todas as regiões e maior ou igual a 62,4 \% em todos os quartos de desempenho ).
Quanto aos enunciados das questões do Componente de Formação Geral ( Questão 4 ), as opiniões foram positivas, já que 76,8 \% dos alunos de Bacharelado avaliados consideraram todos ou a maioria dos enunciados das questões “ claros e objetivos ” ( Gráfico 4.13b, Gráfico 4.14b, e, no Anexo IIB, a Tabela II.7b ).
Na análise regional, a percentagem de estudantes de Bacharelado que avaliaram que todos ou a maioria dos enunciados das questões do Componente de Formação Geral estavam “ claros e objetivos ” variou de 76,1 \%, na região Sul, a 80,0 \%, na região Centro-Oeste.
A análise das percepções dos estudantes de Bacharelado sobre a clareza e a objetividade dos enunciados permite afirmar que todos ou a maioria dos enunciados de questões relativas ao Componente de Formação Geral foram considerados “ claros e objetivos ” entre a maior parte dos respondentes ( maior ou igual a 76,1 \%, em todas as regiões, e maior ou igual a 70,4 \% em todos os quartos de desempenho ).
Segundo o desempenho, observa -se que a proporção dos que emitiram essa opinião cresce, conforme o desempenho aumenta, com diferenças estatisticamente significativas entre o quarto superior de desempenho e os demais.
No quarto superior, a clareza e a objetividade de todos ou da maioria dos enunciados das questões foram percebidas por 81,8 \%.
Segundo o desempenho, observa -se que a proporção dos que emitiram essa opinião cresce conforme o desempenho aumenta, com diferenças estatisticamente significativas entre o quarto inferior e os demais.
No quarto superior, a clareza e a objetividade de todos ou da maioria dos enunciados das questões foram percebidas por 79,7 \%.
A proporção dos presentes de Licenciatura que consideraram que todos ou a maioria dos enunciados das questões do Componente de Formação Geral estavam “ claros e objetivos ” foi maior entre os estudantes de Universidades, chegando a 72,9 \%, enquanto entre os de Faculdades, a Organização Acadêmica com a menor incidência, esse índice atingiu 50,0 \%.
No Gráfico 4.15a, é possível observar que as diferenças entre todos os tipos de Organização Acadêmica não são estatisticamente significativas ( ver também Tabela II.8a, no Anexo IIA ).
A proporção dos presentes de Bacharelado que consideraram que todos ou a maioria dos enunciados das questões do Componente de Formação Geral estavam “ claros e objetivos ” foi maior entre os estudantes de Centros Universitários, atingindo a proporção de 79,9 \%, enquanto entre os de Faculdades, a Organização Acadêmica com a menor incidência, a proporção foi de 72,5 \%.
No Gráfico 4.15b, é possível observar que há diferença estatisticamente significativa entre os Centros Universitários e as Universidades, e entre os Centros Universitários e as Faculdades ( ver também Tabela II.8b, no Anexo IIB ).
O percentual de estudantes de Licenciatura que consideraram que todos ou a maioria dos enunciados das questões do Componente de Formação Geral estavam “ claros e objetivos ” foi maior entre os estudantes de cursos de IES Privadas ( 72,1 \% ), mas sem diferença estatisticamente significativa da proporção alcançada entre os de IES Públicas ( 71,1 \% ).
Um maior detalhamento pode ser encontrado no Gráfico 4.16a e na Tabela II.8a, no Anexo IIA.
O percentual de estudantes de Bacharelado que consideraram que todos ou a maioria dos enunciados das questões do Componente de Formação Geral estavam “ claros e objetivos ” foi ligeiramente maior entre os estudantes de cursos de IES Privadas ( 76,9 \% ), sem diferença estatisticamente significativa da proporção registrada entre os de IES Públicas ( 76,6 \% ).
Um maior detalhamento pode ser encontrado no Gráfico 4.16b e na Tabela II.8b, no Anexo IIB. 4.3.2 COMPONENTE DE CONHECIMENTO ESPECÍFICO Em relação aos enunciados das questões do Componente de Conhecimento Específico da prova, para 68,3 \% dos estudantes de Licenciatura avaliados da Área de Ciência da Computação, a clareza e a objetividade ( Questão 5 ) estavam presentes em todas ou na maioria das questões ( Gráfico 4.17a a Gráfico 4.20a, e no Anexo IIA, as Tabelas II.9a e II.10a ).
A maioria dos estudantes de Licenciatura de todas as Grandes Regiões brasileiras consideraram “ claros e objetivos ” todos ou a maioria dos enunciados das questões do Componente de Conhecimento Específico da prova, percentual sempre maior ou igual a 66,6 \%.
Não há diferença estatisticamente significativa entre as regiões.
Em relação aos enunciados das questões do Componente de Conhecimento Específico da prova, para 79,2 \% dos estudantes de Bacharelado avaliados da Área de Ciência da Computação, a clareza e a objetividade ( Questão 5 ) estavam presentes em todas ou na maioria das questões ( Gráfico 4.17b a Gráfico 4.20b, e no Anexo IIB, as Tabelas II.9b e II.10b ).
A maioria dos estudantes de Bacharelado de todas as Grandes Regiões brasileiras consideraram “ claros e objetivos ” todos ou a maioria dos enunciados das questões do Componente de Conhecimento Específico da prova, percentual sempre maior ou igual a 76,4 \%.
A diferença entre a região Norte e a região Centro-Oeste é estatisticamente significativa, assim como entre a região Sudeste e as regiões Nordeste, Sul e Centro-Oeste.
A proporção de estudantes de Licenciatura que consideraram os enunciados das questões “ claros e objetivos ” apresenta uma tendência crescente em relação ao aumento de desempenho : mais elevada no quarto superior ( 78,6 \% ), se comparada ao quarto inferior de desempenho ( 62,3 \% ).
As diferenças entre o quarto superior e os demais quartos de desempenho são estatisticamente significativas.
A proporção de estudantes de Bacharelado que consideraram os enunciados das questões “ claros e objetivos ” apresenta uma tendência crescente em relação ao aumento de desempenho : mais elevada no quarto superior ( 89,3 \% ), quando comparada ao quarto inferior de desempenho ( 70,3 \% ).
As diferenças entre todos os quartos de desempenho são estatisticamente significativas.
A proporção dos presentes de Licenciatura que consideraram que todos ou a maioria dos enunciados das questões do Componente de Conhecimento Específico estavam “ claros e objetivos ” foi maior entre os estudantes de Universidades, atingindo 71,5 \%.
Entre os estudantes de Faculdades, a Organização Acadêmica com a menor incidência, essa proporção atingiu 50,0 \%.
No Gráfico 4.19a, é possível observar que as diferenças entre todos os tipos de Organização Acadêmica não são estatisticamente significativas ( ver também Tabela II.10a, no Anexo IIA ).
A proporção dos presentes de Bacharelado que consideraram que todos ou a maioria dos enunciados das questões do Componente de Conhecimento Específico estavam “ claros e objetivos ” foi maior entre os estudantes de Universidades, tendo alcançado 80,3 \%, enquanto entre os de Faculdades, a Organização Acadêmica com a menor incidência, foi de 70,7 \%.
No Gráfico 4.19b, é possível observar que há diferença estatisticamente significativa entre as Universidades, os Centros Universitários e as Faculdades ; entre os Centros Universitários e as Faculdades ; e entre as Faculdades e os CEFET/IF ( ver também Tabela II.10b, no Anexo IIB ).
O percentual de estudantes de Licenciatura que consideraram que todos ou a maioria dos enunciados das questões do Componente de Conhecimento Específico estavam “ claros e objetivos ” foi maior entre os estudantes de cursos de IES Privadas ( 70,7 \% ), sem diferença estatisticamente significativa da proporção registrada em IES Públicas ( 68,0 \% ).
Ver Gráfico 4.20a e Tabela II.10a, no Anexo IIA, para maior detalhamento.
O percentual de estudantes de Bacharelado que consideraram que todos ou a maioria dos enunciados das questões do Componente de Conhecimento Específico estavam “ claros e objetivos ” foi maior entre os estudantes de cursos de IES Públicas ( 83,7 \% ), uma diferença estatisticamente significativa da proporção registrada entre os de IES Privadas ( 75,4 \% ).
Ver Gráfico 4.20b e Tabela II.10b, no Anexo IIB, para um maior detalhamento. 4.4 SUFICIÊNCIA DAS INFORMAÇÕES/INSTRUÇÕES FORNECIDAS Ao se avaliarem as informações/instruções fornecidas para a resolução das questões ( Questão 6 ), constatou -se que 74,3 \% dos respondentes de Licenciatura da Área de Ciência da Computação de todo o Brasil afirmaram que estas eram até excessivas ou suficientes em todas ou na maioria das questões ( Gráfico 4.21a, Gráfico 4.22a e, no Anexo IIA, a Tabela II.11a ).
Quanto à distribuição de respondentes pelas Grandes Regiões, observa -se que a proporção de estudantes que consideraram as informações/instruções fornecidas até excessivas ou suficientes em todas ou na maioria das questões foi sempre superior ou igual a 68,3 \%, chegando a 86,9 \%, na região Centro-Oeste.
Existem diferenças estatisticamente significativas apenas entre a região Centro-Oeste e as regiões Norte, Nordeste e Sudeste.
Ao avaliarem as informações/instruções fornecidas para a resolução das questões ( Questão 6 ), 83,7 \% dos respondentes de Bacharelado da Área de Ciência da Computação de todo o Brasil afirmaram que estas eram até excessivas ou suficientes em todas ou na maioria das questões ( Gráfico 4.21b, Gráfico 4.22b, e, no Anexo IIB, a Tabela II.11b ).
Quanto à distribuição de respondentes pelas Grandes Regiões, observa -se que a proporção de estudantes que consideraram as informações/instruções fornecidas até excessivas ou suficientes em todas ou na maioria das questões foi sempre superior ou igual a 81,9 \%, chegando a 87,7 \%, na região Sul.
Existem diferenças estatisticamente significativas apenas entre a região Sudeste e a regiões Nordeste e Sul.
Levando -se em conta o desempenho dos participantes de Licenciatura, observa -se uma tendência crescente das proporções de participantes que avaliaram as informações/instruções como até excessivas ou suficientes em todas ou na maioria das questões, havendo diferença estatisticamente significativa apenas entre o primeiro e o quarto superior de desempenho, conforme mostrado no Gráfico 4.22a.
O percentual foi mais elevado no quarto superior ( 80,1 \% ), percentual superior à média nacional ( 74,3 \% ).
Já no quarto inferior, a suficiência das informações/instruções declarada como até excessiva, em todas ou na maioria das questões foi percebida por 68,1 \% dos respondentes.
Levando -se em conta o desempenho dos participantes de Bacharelado, observa -se uma tendência crescente das proporções de participantes que avaliaram as informações/instruções como até excessivas ou suficientes em todas ou na maioria das questões, havendo diferenças estatisticamente significativas entre o primeiro e os demais quartos de desempenho, assim como entre o quarto superior e os segundo e terceiro quartos, como mostrado no Gráfico 4.22b.
O percentual foi mais elevado no quarto superior ( 90,0 \% ), percentual superior à média nacional ( 83,7 \% ).
Já no quarto inferior, a suficiência das informações/instruções declarada como até excessiva, em todas ou na maioria das questões foi percebida por 77,6 \% dos respondentes.
Considerando -se as informações/instruções fornecidas para a resolução das questões ( Questão 6 ), a proporção de estudantes que consideraram que estas eram até excessivas ou suficientes em todas ou na maioria das questões foi maior entre os estudantes de Licenciatura de Centros Universitários, totalizando 78,2 \%.
As Faculdades foram a Organização Acadêmica em que houve a menor proporção : 50,0 \%.
No Gráfico 4.23a, é possível observar que apenas a diferença entre Universidades ( 77,5 \% ) e CEFET/IF ( 70,4 \% ) é estatisticamente significativa ( ver também Tabela II.12a, no Anexo IIA ).
Considerando -se as informações/instruções fornecidas para a resolução das questões ( Questão 6 ), a proporção de estudantes que consideraram que estas eram até excessivas ou suficientes em todas ou na maioria das questões foi maior entre os estudantes de Bacharelado de Universidades : 84,7 \% ; já entre os de Faculdades, a Organização Acadêmica com a menor incidência, o índice atingiu 78,0 \%.
No Gráfico 4.23b, é possível observar que somente existe diferença estatisticamente significativa entre as Universidades e os Centros Universitários ( 81,5 \% ) e as Faculdades ( ver também Tabela II.12b, no Anexo IIB ).
O percentual de estudantes de Licenciatura que consideraram que as informações/instruções fornecidas para a resolução das questões eram até excessivas ou suficientes em todas ou na maioria das questões foi maior entre os estudantes de cursos de IES Privadas ( 77,6 \% ), não havendo diferença estatisticamente significativa da proporção registrada em IES Públicas ( 73,9 \% ) ( Ver Gráfico 4.24a e Tabela II.12a, no Anexo IIA ).
O percentual de estudantes de Bacharelado que consideraram que as informações/instruções fornecidas para a resolução das questões eram até excessivas ou suficientes em todas ou na maioria das questões foi maior entre os estudantes de cursos de IES Públicas ( 87,3 \% ), uma diferença estatisticamente significativa da proporção registrada entre os de IES Privadas ( 80,7 \% ) ( ver Gráfico 4.24b e Tabela II.12b, no Anexo IIB ). 4.5 DIFICULDADE ENCONTRADA AO RESPONDER À PROVA Indagados sobre as dificuldades com as quais se depararam ao responder à prova ( Questão 7 ), 22,7 \% dos estudantes de Licenciatura apontaram o Desconhecimento do conteúdo.
Para 50,4 \%, a Forma diferente de abordagem do conteúdo foi indicada como dificuldade.
Já a Falta de motivação para fazer a prova foi a dificuldade apontada por 17,2 \% dos respondentes.
Considerando -se todo o Brasil, 6,9 \% dos respondentes afirmaram que não tiveram qualquer tipo de dificuldade para responder à prova ( Tabelas II.13a e II.14a, no Anexo IIA ).
Nos Gráficos 4.25a e 4.28a, são apresentados os percentuais de estudantes de Licenciatura que apontaram o Desconhecimento do conteúdo como “ dificuldade ao responder à prova ”.
Na análise por Grandes Regiões, o percentual de inscritos e presentes de Licenciatura que apontaram o Desconhecimento do conteúdo como “ dificuldade ao responder à prova ” não superou 27,4 \%.
Os percentuais variaram de 19,9 \%, na região Nordeste, a 27,4 \%, na região Sul.
A Forma diferente de abordagem do conteúdo foi a escolha modal dos estudantes de Licenciatura, com percentuais que variaram de 37,3 \% ( região Centro-Oeste ) a 56,7 \% ( região Norte ).
O percentual de alunos que citaram a Falta de motivação como dificuldade variou de 12,6 \% ( região Sul ) a 26,5 \% ( região Centro-Oeste ).
Os que declararam não ter qualquer tipo de dificuldade para responder à prova variaram de 2,8 \%, na região Norte, a 8,0 \%, na região Sudeste.
Indagados sobre as dificuldades com as quais se depararam ao responder à prova ( Questão 7 ), 23,9 \% dos estudantes de Bacharelado apontaram o Desconhecimento do conteúdo.
Para 36,2 \%, a Forma diferente de abordagem do conteúdo foi indicada como dificuldade.
Já a Falta de motivação para fazer a prova foi a dificuldade apontada por 24,5 \% dos respondentes.
Considerando -se todo o Brasil, 12,2 \% dos respondentes afirmaram que não tiveram qualquer tipo de dificuldade para responder à prova ( Tabelas II.13b e II.14b, no Anexo IIB ).
Nos Gráficos 4.25b e 4.28b, são apresentados os percentuais de estudantes de Bacharelado que apontaram o Desconhecimento do conteúdo como “ dificuldade ao responder à prova ”.
Na análise por Grandes Regiões, o percentual de inscritos e presentes de Bacharelado que apontaram o Desconhecimento do conteúdo como “ dificuldade ao responder à prova ” não superou 28,1 \%.
Os percentuais variaram de 20,2 \%, na região Sul, a 28,1 \%, na região Norte.
A Forma diferente de abordagem do conteúdo foi a escolha modal dos estudantes de Bacharelado, com percentuais que variaram de 31,0 \% ( região Nordeste ) a 39,9 \% ( região Norte ).
O percentual de alunos que citaram a Falta de motivação como dificuldade variou de 21,8 \% ( região Sudeste ) a 29,6 \% ( região Nordeste ).
Os que declararam não ter qualquer tipo de dificuldade para responder à prova variaram de 7,6 \%, na região Norte a 13,3 \%, na região Nordeste.
Em relação aos quartos de desempenho, o Desconhecimento do conteúdo foi a opção escolhida por 18,5 \% dos estudantes de Licenciatura do quarto inferior e por 28,0 \% do quarto superior, caracterizando uma tendência crescente.
É observada diferença estatisticamente significativa apenas entre o quarto inferior e o quarto superior de desempenho.
A alternativa modal entre os alunos de Licenciatura, quando agregados pelos quartos de desempenho, para a dificuldade encontrada, foi causada pela Forma diferente de abordagem do conteúdo : variando de 47,3 \%, no terceiro quarto, a 53,6 \%, no segundo quarto.
Em relação aos quartos de desempenho, o Desconhecimento do conteúdo foi a opção escolhida por 22,3 \% dos estudantes de Bacharelado do quarto inferior e por 26,5 \% do quarto superior, caracterizando uma tendência crescente.
É observada diferença estatisticamente significativa apenas entre o quarto superior e os dois quartos inferiores de desempenho.
A alternativa modal entre os alunos de Bacharelado, quando agregados pelos quartos de desempenho, para a dificuldade encontrada, foi causada pela Forma diferente de abordagem do conteúdo : variando de 30,8 \%, no quarto superior, a 40,1 \%, no primeiro quarto.
Na análise por Organização Acadêmica, o percentual de inscritos e presentes de Licenciatura que apontaram o Desconhecimento do conteúdo como “ dificuldade ao responder à prova ” não superou 25,7 \%.
Os percentuais variaram de 0,0 \%, nas Faculdades, a 25,7 \%, nos Centros Universitários.
A Forma diferente de abordagem do conteúdo foi a escolha modal dos estudantes de Licenciatura, com percentuais que variaram de 42,9 \% ( Centros Universitários ) a 100,0 \% ( Faculdades ).
O percentual de alunos de Licenciatura que citaram a Falta de motivação como dificuldade variou de 0,0 \% ( Faculdades ) a 17,8 \% ( Universidades ).
Os que declararam não ter qualquer tipo de dificuldade para responder à prova variaram de 0,0 \%, nas Faculdades, a 7,5 \%, nas Universidades.
No Gráfico 4.27a, é possível observar que há diferenças estatisticamente significativas entre as Faculdades e os demais tipos de Organização Acadêmica ( ver também Tabela II.14a, no Anexo IIA ).
Na análise por Organização Acadêmica, o percentual de inscritos e presentes de Bacharelado que apontaram o Desconhecimento do conteúdo como “ dificuldade ao responder à prova ” não superou 25,8 \%.
Os percentuais variaram de 22,4 \%, para CEFET/IF a 25,8 \%, nas Faculdades.
A Forma diferente de abordagem do conteúdo foi a escolha modal dos estudantes de Bacharelado, com percentuais que variaram de 34,6 \% ( Universidades ) a 42,8 \% ( Faculdades ).
O percentual de alunos de Bacharelado que citaram a Falta de motivação como dificuldade variou de 18,1 \% ( Faculdades ) a 26,1 \% ( Universidades ).
Os que declararam não ter qualquer tipo de dificuldade para responder à prova variaram de 9,6 \%, nos CEFET/IF a 12,5 \%, nos Universidades.
No Gráfico 4.27b, é possível observar que não há diferenças estatisticamente significativas entre os tipos de Organização Acadêmica ( ver também Tabela II.14b, no Anexo IIB ).
Em relação às Categorias Administrativas, o Desconhecimento do conteúdo foi a opção escolhida por 25,5 \% dos estudantes de Licenciatura nas IES Privadas e por 22,3 \% nas IES Públicas.
Não são observadas diferenças estatisticamente significativas entre essas duas categorias.
A alternativa modal para os alunos de Licenciatura, quando agregados pelas Categorias Administrativas, para a dificuldade encontrada, foi causada pela Forma diferente de abordagem do conteúdo : 51,2 \%, entre os estudantes de IES Públicas, e 43,4 \%, entre os de IES Privadas ( ver Gráfico 4.28a e Tabela II.14a, no Anexo IIA ).
Em relação às Categorias Administrativas, o Desconhecimento do conteúdo foi a opção escolhida por 24,4 \% dos estudantes de Bacharelado, nas IES Privadas, e por 23,2 \%, nas IES Públicas.
Não são observadas diferenças estatisticamente significativas entre essas duas categorias.
A alternativa modal para os alunos de Bacharelado, quando agregados pelas Categorias Administrativas, para a dificuldade encontrada, foi causada pela Forma diferente de abordagem do conteúdo : 30,8 \%, entre os estudantes de IES Públicas, e 40,8 \%, entre os de IES Privadas ( ver Gráfico 4.28b e Tabela II.14b, no Anexo IIB ). 4.6 CONTEÚDOS DAS QUESTÕES OBJETIVAS DA PROVA Ao se analisarem os conteúdos das questões objetivas da prova ( Questão 8 ), verifica -se que um percentual muito pequeno dos estudantes de Licenciatura avaliados, apenas 11,2 \%, afirmaram que não estudaram ainda a maioria desses conteúdos ( Gráficos 4.29a a Gráfico 4.32a, e as Tabelas II.15a e II.16a, no Anexo IIA ).
A maioria ( 43,3 \% ) afirmou ter estudado e aprendido muitos ou todos os conteúdos avaliados.
Na análise por Grande Região, a proporção de respondentes de Licenciatura que escolheram a opção não estudou ainda a maioria desses conteúdos foi pequena.
Observa -se que, nas regiões Norte ( 17,3 \% ) e Sudeste ( 11,7 \% ), as proporções foram maiores do que a da média nacional ( 11,2 \% ).
Nas regiões Nordeste ( 10,9 \% ), Sul ( 7,6 \% ) e Centro-Oeste ( 6,0 \% ), a situação foi inversa, com proporções mais baixas.
São observadas diferenças estatisticamente significativas apenas entre a região Norte e as regiões Sul e Centro-Oeste.
Em todas as regiões, a maioria dos presentes afirmou ter estudado e aprendido muitos ou todos os conteúdos, com proporções variando entre 30,7 \%, na região Norte, a 60,7 \%, na região Sul.
Na análise dos conteúdos das questões objetivas da prova ( Questão 8 ), um percentual muito pequeno dos estudantes de Bacharelado avaliados, apenas 3,1 \%, afirmaram que não estudaram ainda a maioria desses conteúdos ( Gráficos 4.29b a Gráfico 4.32b, e as Tabelas II.15b e II.16b, no Anexo IIB ).
A maioria ( 63,7 \% ) afirmou ter estudado e aprendido muitos ou todos os conteúdos avaliados.
Na análise por Grande Região, a proporção de respondentes de Bacharelado que escolheram a opção não estudou ainda a maioria desses conteúdos foi pequena.
Observa -se que, nas regiões Nordeste ( 2,0 \% ), Sul ( 2,2 \% ) e Centro-Oeste ( 2,6 \% ), as proporções foram menores do que a da média nacional ( 3,1 \% ).
Nas regiões Norte e Sudeste, a situação foi inversa : uma proporção mais alta que a da média nacional ( 3,6 \% e 3,8 \%, respectivamente ).
São observadas diferenças estatisticamente significativas entre a região Sudeste e as regiões Nordeste e Sul.
Em todas as regiões, a maioria dos presentes afirmou ter estudado e aprendido muitos ou todos os conteúdos, com proporções variando entre 59,3 \%, na região Norte, a 73,5 \%, na região Sul.
Considerando -se, separadamente, as opiniões de estudantes de Licenciatura dos quatro quartos de desempenho, observa -se que, no quarto inferior, 16,1 \% ofereceram como resposta que não estudaram ainda a maioria desses conteúdos, sendo 6,0 \% os do quarto superior com a mesma resposta.
Apenas a diferença registrada entre os alunos que optaram por esse motivo de dificuldade no quarto superior em relação aos dois quartos inferiores de desempenho é estatisticamente significativa.
Tendo -se em conta o quarto superior, 56,0 \% dos alunos de Licenciatura afirmaram ter estudado e aprendido muitos ou todos os conteúdos.
No outro extremo, no primeiro quarto, 36,8 \% optaram pelas mesmas categorias.
As proporções são crescentes com o desempenho.
Considerando -se, separadamente, as opiniões de estudantes de Bacharelado dos quatro quartos de desempenho, observa -se que, no quarto inferior, 6,4 \% ofereceram como resposta a alternativa não estudou ainda a maioria desses conteúdos, sendo 0,7 \% os do quarto superior com a mesma resposta.
As diferenças entre os quartos de desempenho são estatisticamente significativas.
Tendo -se em conta o quarto superior, 82,0 \% dos alunos de Bacharelado afirmaram ter estudado e aprendido muitos ou todos os conteúdos.
No outro extremo, no primeiro quarto, 47,2 \% optaram pelas mesmas categorias.
As proporções são crescentes com o desempenho.
Na análise por Organização Acadêmica, a proporção de respondentes de Licenciatura que escolheu a opção não estudou ainda a maioria desses conteúdos foi pequena.
Observa -se que as proporções nas Universidades ( 8,6 \% ) e Faculdades ( 0,0 \% ) foram menores do que a da média nacional ( 11,2 \% ).
Nos Centros Universitários ( 19,4 \% ) e nos CEFET/IF ( 11,9 \% ), a situação foi inversa, com proporção mais alta.
No Gráfico 4.31a, é possível observar que as diferenças entre Faculdades e os demais tipos de Organização Acadêmica são estatisticamente significativas, assim como entre Universidades e Centros Universitários.
Em todos os tipos de Organização Acadêmica, a maioria dos presentes afirmou ter estudado e aprendido muitos ou todos os conteúdos, com proporções variando entre 37,7 \%, nos CEFET/IF, a 100,0 \%, nas Faculdades. ( Ver também Tabela II.16a, no Anexo IIA ).
Na análise por Organização Acadêmica, a proporção de respondentes de Bacharelado que escolheram a opção não estudou ainda a maioria desses conteúdos foi pequena.
Observa -se que, nas Faculdades ( 7,0 \% ) e nos Centros Universitários ( 3,4 \% ), as proporções foram maiores do que a da média nacional ( 3,1 \% ).
Nas Universidades ( 2,8 \% ) e nos CEFE/IF ( 1,5 \% ), a situação foi inversa : as proporções foram mais baixas.
No Gráfico 4.31b, é possível observar que a diferença entre as Faculdades e os demais tipos de Organização Acadêmica é estatisticamente significativa ( ver também Tabela II.16b, no Anexo IIB ).
Em todos os tipos de Organização Acadêmica, a maioria dos presentes afirmou ter estudado e aprendido muitos ou todos os conteúdos, com proporções variando entre 49,5 \%, nas Faculdades, a 66,1 \%, nos Universidades.
Em relação às Categorias Administrativas, observa -se que, nas IES Públicas, 10,3 \% ofereceram como resposta a alternativa não estudou ainda a maioria desses conteúdos, sendo 18,8 \% nas IES Privadas com a mesma resposta.
A diferença entre os alunos de Licenciatura que optaram por esse motivo de dificuldade nas duas categorias é estatisticamente significativa.
Considerando -se as IES Públicas, 43,4 \% dos alunos de Licenciatura afirmaram ter estudado e aprendido muitos ou todos os conteúdos, sendo de 43,1 \% a proporção dos que escolheram essa resposta nas IES Privadas ( ver Gráfico 4.32a e Tabela II.16a, no Anexo IIA ).
Em relação às Categorias Administrativas, observa -se que, nas IES Públicas, 1,5 \% dos estudantes ofereceram como resposta a alternativa não estudou ainda a maioria desses conteúdos, sendo 4,5 \% nas IES Privadas com a mesma resposta.
A diferença entre os alunos de Bacharelado que optaram por esse motivo de dificuldade nas duas categorias é estatisticamente significativa.
Considerando -se as IES Públicas, 72,9 \% dos alunos de Bacharelado afirmaram ter estudado e aprendido muitos ou todos os conteúdos.
Na outra categoria, 55,9 \% optaram pelas mesmas categorias ( ver Gráfico 4.32b e Tabela II.16b, no Anexo IIB ). 4.7 TEMPO GASTO PARA CONCLUIR A PROVA Ao responderem sobre o tempo de conclusão da prova ( Questão 9 ), mais de três quartos dos estudantes de Licenciatura ( 76,0 \% ) afirmaram ter gasto Entre duas e quatro horas ( Gráfico 4.33a a Gráfico 4.36a e, no Anexo IIA, as Tabelas II.17a e II.18a ).
Considerando -se as cinco Grandes Regiões brasileiras, os percentuais dos que utilizaram Entre duas e quatro horas para finalizar a prova nas regiões Sudeste ( 75,1 \% ) e Centro-Oeste ( 66,3 \% ) foram inferiores ao percentual nacional.
Nas demais Grandes Regiões, os percentuais de alunos de Licenciatura que dispensaram Entre duas e quatro horas para concluir a prova ficaram acima do nacional : 80,1 \% na região Norte ; 76,2 \% na região Nordeste ; e 78,0 \%, na região Sul, como mostrado no Gráfico 4.33a.
Não há diferenças estatisticamente significativas entre as regiões.
Ao responderem sobre o tempo de conclusão da prova ( Questão 9 ), mais de dois terços dos estudantes de Bacharelado ( 70,8 \% ) afirmaram ter gasto Entre duas e quatro horas ( Gráfico 4.33b a Gráfico 4.36b e, no Anexo IIB, as Tabelas II.17b e II.18b ).
Considerando -se as cinco Grandes Regiões brasileiras, os percentuais dos que utilizaram Entre duas e quatro horas para finalizar a prova na região Sudeste ( 69,0 \% ) foram inferiores ao percentual nacional.
Nas demais Grandes Regiões, o percentual de alunos de Bacharelado que dispensaram Entre duas e quatro horas para concluir a prova ficou acima : 73,5 \% na região Norte, 72,7 \% na região Nordeste ; 72,1 \%, na região Sul e 75,2 \% na região Centro-Oeste, como mostrado no Gráfico 4.33b.
A diferença entre a região Sudeste e as regiões Nordeste e Centro-Oeste é estatisticamente significativa.
Uma vez analisadas as alternativas escolhidas pelos estudantes de Licenciatura que se situam nos diferentes quartos de desempenho, observa -se uma tendência crescente : uma proporção maior de participantes no quarto superior declarou ter gasto Entre duas e quatro horas para concluir a prova quando comparada à proporção de participantes nos quartos inferiores.
Há diferenças estatisticamente significativas entre o quarto superior e os dois primeiros quartos de desempenho.
Uma vez analisadas as alternativas escolhidas pelos estudantes de Bacharelado que se situam nos diferentes quartos de desempenho, observa -se uma tendência crescente : uma proporção maior de participantes no quarto superior declarou ter gasto Entre duas e quatro horas para concluir a prova quando comparada à proporção de participantes nos quartos inferiores.
Há diferenças estatisticamente significativas entre o quarto inferior e os demais quartos de desempenho, assim como entre o quarto superior e os demais quartos de desempenho.
Na análise por Organização Acadêmica, os percentuais dos que utilizaram Entre duas e quatro horas para finalizar a prova nos CEFET/IF ( 78,5 \% ) foi superior ao percentual nacional ( 76,0 \% ).
Nas demais Organizações Acadêmicas, o percentual foi inferior : Universidades ( 75,9 \% ), Centros Universitários ( 65,7 \% ) e Faculdades ( 50,0 \% ).
Como mostrado no Gráfico 4.35a, apenas há diferença estatisticamente significativa entre as proporções dos Centros Universitários e as dos CEFET/IF.
Na análise por Organização Acadêmica, os percentuais dos que utilizaram Entre duas e quatro horas para finalizar a prova nas Universidades ( 70,7 \% ) e Faculdades ( 68,7 \% ) foram inferiores ao percentual nacional ( 70,8 \% ).
Nos demais tipos de Organização Acadêmica, o percentual de alunos de Bacharelado que dispensaram Entre duas e quatro horas para concluir a prova alcançou 71,0 \%, nos Centros Universitários e 74,3 \%, nos CEFET/IF, conforme mostrado no Gráfico 4.35b.
As diferenças entre os tipos de Organizações Acadêmicas não são estatisticamente significativas.
Em relação às Categorias Administrativas, observa -se que, nas IES Públicas, 77,2 \% declararam ter gasto Entre duas e quatro horas para concluir a prova, enquanto 66,2 \% nas IES Privadas ofereceram a mesma resposta.
A diferença entre os alunos de Licenciatura que optaram por esse motivo de dificuldade nas duas categorias é estatisticamente significativa ( ver Gráfico 4.36a e Tabela II.18a, no Anexo IIA ).
Em relação às Categorias Administrativas, observa -se que, nas IES Públicas, 72,1 \% declararam ter gasto Entre duas e quatro horas para concluir a prova, enquanto 69,8 \% nas IES Privadas forneceram a mesma resposta.
A diferença entre os alunos de Bacharelado que optaram por esse motivo de dificuldade nas duas categorias não é estatisticamente significativa. ( Ver Gráfico 4.36b e Tabela II.18b, no Anexo IIB ).

% matched lemmas: repor
\end{textsample}
